\section{Introduction}
\label{sec:intro}

The \nusquids\ package~\cite{nusquids2021} solves the neutrino
density-matrix evolution equation in the interaction picture for
arbitrary trajectories through dense media, supporting coherent
oscillations in vacuum and matter, neutral- and charged-current
absorption, neutral-current cascading, tau regeneration, and the
Glashow resonance for electron antineutrinos. The CPU implementation
relies on the \squids~\cite{squids} library, which provides an
adaptive Runge--Kutta--Fehlberg (RKF45) integrator built on
\textsc{GSL}~\cite{gsl}. For the multi-energy, multi-zenith grids that
underlie IceCube atmospheric oscillation analyses, the wall-time of a
typical computation scales as $\mathcal{O}(n_{c_\theta}\, n_E\, n_\rho\,
n_\nu)$ density-matrix evolutions, which becomes a meaningful fraction
of the analysis turnaround as grid resolution grows.

This article describes a GPU port of the \nusquids\ solver. The
implementation maps each independent zenith trajectory to a CUDA
thread block, parallelises the energy and flavor dimensions across
threads within the block, and preserves the adaptive step size
control of the CPU integrator via a Richardson extrapolation
step-doubling scheme on the device. We pay particular attention to
the cascade source terms (NC, tau decay, Glashow), which couple
energy bins and therefore require cooperative reductions; we
demonstrate that refreshing these source terms at every Runge--Kutta
substage --- mirroring the per-evaluation \texttt{PreDerive} update
on the CPU --- is necessary for agreement at the sub-percent level on
tau-regeneration-dominated scenarios.

A first GPU port of \nusquids\ -- \texttt{CUDA}$\nu$\texttt{SQuIDS} --
was published by Kallenborn~\textit{et al.}~\cite{kallenborn2019},
who reported up to $71\times$ speedup on a Titan V for low-energy
oscillation-only computations and $21$--$24\times$ speedup for
high-energy computations including non-coherent interactions on a
$200\times200$ grid. The present work builds on that codebase
(LGPLv3) and revisits the integrator structure to obtain a higher
speedup specifically in the interactions-included regime, which is
the dominant cost in IceCube atmospheric oscillation analyses.
We compare to their published numbers in
Sec.~\ref{sec:perf_priorwork}.

The remainder of the paper is organised as follows.
Section~\ref{sec:gpu_impl} describes the GPU architecture and the
substage-refresh integrator. Section~\ref{sec:validation} validates
the implementation against the CPU reference on five representative
scenarios. Section~\ref{sec:performance} quantifies the speedup on a
range of NVIDIA architectures (A100 MIG and full, % TODO: H100, H200).
Section~\ref{sec:conclusion} concludes.
